% GENERATED by tools/failure_weight_decomposition_figure.py from
%   data/ogasp/controller/outcome_rules.json + lifecycle_consensus.json
%   through mtdsim.l3_simulation.controller.rules (overlay version v4_failure_only).
% Do not hand-edit; regenerate. Requires booktabs + graphicx (both in the preamble).

\begin{table}[htbp]
\centering
\caption[The rules of the failure matrix]{The rules that set the failure matrix (Section~\ref{subsec:failure-matrix}). After a failure at one tactic, the first rule that matches the move to the next tactic gives its value, which the distance factor of Table~\ref{tab:overlay-distance-kernel} then scales. Back, within and forward are by lifecycle stage. The letters are those of Figure~\ref{fig:failure-weight-decomposition}(a).}
\label{tab:overlay-failure-rules}
\tablestyle
\begin{tabular}{@{}c P{0.40\textwidth} r P{0.40\textwidth}@{}}
\toprule
Rule & Applies to & Value & Reason \\
\midrule
A & After a failed initial access, a move to a tactic that needs a foothold & 0.02 & No foothold was gained, so these moves are almost out of reach. \\
B & After a failed reconnaissance, a move to initial access & 0.4 & Nothing was found, so breaking in is harder but still possible. \\
C & After a failed reconnaissance, a move to a tactic that needs a foothold & 0.05 & With no foothold, a move deep into the network is implausible. \\
D & After a failure with a foothold, a move back to a tactic without one & 0.25 & Falling right back to before the foothold is an uncommon response. \\
E & Any other move back to execution & 0.35 & Running code again on a held foothold continues the attack. \\
F & Any other move back: from initial access to preparation, or from actions on objectives to post-intrusion operations & 0.9 & Going back one stage is the common response to a failure. \\
G & A move within the same stage & 0.7 & Another tactic of the stage is an alternative route. \\
H & A move forward, with a foothold & 0.35 & Advancing straight after a failure is possible but uncommon. \\
I & A move forward, without a foothold & 0.3 & The step the next one depends on did not succeed. \\
\bottomrule
\end{tabular}
\end{table}

% No success-rules table and no success-set table for v4_failure_only:
%   the success table of this version is the multiplicative identity (every
%   cell 1.0, rule `passthrough`) — on a success verdict the token routes on
%   the corpus proportions unchanged, and only the failure table is declared.
%   The retired success rules remain in outcome_rules.json, unconsulted.

\begin{table}[htbp]
\centering
\caption[The distance factor of the failure matrix]{The distance factor that scales each rule value by the lifecycle stages a move crosses. A move within a stage or to the next keeps its rule value; each further stage forward multiplies it by $\gamma$, and each further stage back by $\delta$. A factor below $z$ is zero. Appendix~\ref{app:decay-robustness} moves each parameter across its band.}
\label{tab:overlay-distance-kernel}
\tablestyle
\begin{tabular}{@{}l r l@{}}
\toprule
Parameter & Value & Band \\
\midrule
$\gamma$, forward & 0.25 & \{0.1, 0.5\} \\
$\delta$, back & 0.25 & \{0.1, 0.5\} \\
$z$, floor & 0.1 & \{0, 0.05, 0.1\} \\
\bottomrule
\end{tabular}
\end{table}
